% !TeX program = xelatex
\documentclass[11pt,a4paper]{article}
\usepackage[margin=23mm]{geometry}
\usepackage{fontspec}
\setmainfont{Latin Modern Roman}
\usepackage{amsmath,amssymb,mathtools}
\usepackage{xcolor,enumitem,booktabs,longtable,array,xurl,hyperref}
\definecolor{muted}{HTML}{53636C}
\hypersetup{colorlinks=true,urlcolor=blue,linkcolor=blue}
\setlength{\parindent}{0pt}
\setlength{\parskip}{5pt}
\newcommand{\R}{\mathbb R}
\newcommand{\E}{\mathbb E}
\newcommand{\N}{\mathbb N}
\newcommand{\ind}{\mathbf 1}
\newcommand{\Var}{\operatorname{Var}}
\newcommand{\Cov}{\operatorname{Cov}}
\newcommand{\supp}{\operatorname{supp}}
\DeclareMathOperator*{\argmin}{arg\,min}
\DeclareMathOperator*{\argmax}{arg\,max}
\newcommand{\entrymeta}[3]{{\sffamily\small\color{muted}\textbf{\texttt{#1}}\quad|\quad #2\par #3\par}}
\newcommand{\Problemref}[1]{\texttt{#1}}
\newcommand{\JELref}[2]{\texttt{#2} (JEL \texttt{#1})}
\begin{document}
\section*{Exclusionary contracting}
\textbf{Problem ID:} \texttt{L42-195}\quad \textbf{JEL:} \texttt{L42}\par
% BEGIN ECONOMICS STATEMENT
\label{problem:OP-0399}
{\sffamily\small\color{muted}Original subject group: Industrial organization and competition\par}
\label{definitions:problem:30007681}
\textbf{Audit classification:} Published research agenda. The chapter supplies modeling directions, not an exact published conjecture about every exclusionary contract. Entry deterrence and real-cost savings are separate and can coexist.
\subsection*{Definitions and precise research target}
Consider a market with one upstream incumbent, one possible upstream entrant and finitely many downstream retailers. Retailers purchase inputs under contracts specifying quantity-dependent payments and an exclusivity indicator. Let \(r\in\{0,1\}\) denote a regulatory regime: \(r=1\) permits exclusivity and nonlinear payments, while \(r=0\) prohibits exclusivity but preserves a feasible nonlinear-payment class to be declared for each model. Let \(\theta\) collect consumer demand, production technologies, entrant fixed costs, investment opportunities and each firm's information. Firms bargain, choose investments and decide entry under a declared timing protocol. For each regime, let \(\mathcal E_r(\theta)\) contain equilibria consistent with that protocol; restrict comparisons to parameters for which both sets are nonempty. In equilibrium \(e\), define \(W(e)\) as consumer surplus plus producer profits, with profits measured net of real production, entry and investment expenditure. Determine welfare comparisons between regimes allowing retailers to rematch with suppliers, rather than mechanically preserving incumbent contracts. Separately define entry deterrence as lower entrant participation in regime \(1\) and productive coordination as lower real cost or higher quality under an explicitly stated accounting of investment costs. Characterize when both effects coexist and which dominates welfare under a stated equilibrium-selection rule and an explicit mechanism comparison. The problem is to obtain tractable, estimable conditions within this model class and identify data on contract terms and switching that distinguish those mechanisms. It does not relabel previously solved exclusive-dealing examples as universal unresolved theorems.

\textbf{Definitions, assumptions and mathematical qualifications.} This is a modeling and empirical research agenda, not a claim quantified over every game with these market participants. Each contribution must declare its model or restricted class before drawing conclusions, state the parameter domain, and distinguish conclusions for a constructed example from conclusions for every member of a declared class. Model construction, equilibrium existence, welfare comparison, identification, computation and statistical estimation are distinct deliverables; a contribution must state which it addresses and the conditions under which its conclusions hold.

For each model, index retailers by \(j\in\{1,\ldots,n\}\), with declared \(n\geq1\), and specify quantity spaces and their measurable structures. A contract includes a measurable payment schedule \(T_j:Q_j\to\R\) and an exclusivity indicator; supplier-specific schedules and any dependence on other purchases or histories must be specified when used. Declare the feasible schedule classes, other contractual instruments, investment and price choices, exclusivity obligations and their duration, commitment and enforcement rules, information and beliefs, disagreement outcomes, and opportunities to reject, renegotiate and rematch. Give a complete bargaining/investment/entry timing game and its equilibrium concept. Measurability alone is not a feasibility or equilibrium-existence condition. Specify both regimes' admissible contract sets and enforcement rules, including whether regime \(0\) is exactly the no-exclusivity restriction of regime \(1\), and disclose any other differences. Permission to use exclusivity does not require its use; a ban on explicit exclusivity does not by itself remove exclusionary incentives from nonlinear payments. State how such incentives are treated. The same cost, demand and shock primitives apply in both regimes.

For a declared model and common primitives \(\theta\), give a measurable equilibrium space and a probability selection kernel \(\kappa_r(de\mid\theta)\) concentrated on \(\mathcal E_r(\theta)\). Here \(e\) is an equilibrium strategy profile, including beliefs when required, rather than a realized allocation. Declare the outcome-history space and the probability law \(P_r(dh\mid e,\theta)\) induced by equilibrium behavior and within-equilibrium uncertainty. Specify a common comparison horizon and an activity event \(A\) on that history space, such as entry followed by production during a declared period; the chosen event is part of the model. Let \(w_r(h;\theta)\) and \(c_r(h;\theta)\) denote realized welfare and real resource cost under the accounting conventions below, and set
\[
 W(e)=\int w_r(h;\theta)P_r(dh\mid e,\theta),\qquad
 C(e)=\int c_r(h;\theta)P_r(dh\mid e,\theta)
\]
for \(e\in\mathcal E_r(\theta)\), suppressing regime and primitive arguments only in this notation. Define
\[
 E_r(\theta)=\int P_r(A\mid e,\theta)\kappa_r(de\mid\theta),\qquad
 W_r(\theta)=\int W(e)\kappa_r(de\mid\theta),\qquad
 R_r(\theta)=\int C(e)\kappa_r(de\mid\theta).
\]
Require measurable integrands and absolute integrability of welfare and resource costs under the joint selection and outcome law for each reported finite comparison. Comparisons averaged across primitive heterogeneity must additionally declare a common distribution of \(\theta\) and the corresponding integrability conditions. State whether selection is specified, estimated or varied over a declared admissible family; a robust assertion must state its quantifier over that family. Deterrence means \(E_1(\theta)<E_0(\theta)\) at the stated primitives, or the expressly defined population-averaged counterpart. Report empty equilibrium sets as failure of equilibrium existence for the affected regime and solution concept, separately from an empty feasible set; neither permits a selected-equilibrium welfare comparison.

Define consumer surplus through an explicit monetary normalization of consumer utility, with the treatment of income effects stated. Include the incumbent, entrant and all retailers in producer profits, count real production, entry and investment expenditures once, and cancel transfers internal to the market consistently. Specify the numeraire, horizon and any discounting. A productivity comparison must declare its output and quality bundle, the relevant production or investment counterfactual and the resource expenditures included. Lower total cost caused by lower output is not automatically greater efficiency; a cost saving relative to the incumbent's own benchmark is not automatically a reduction in total regime cost. A quality comparison must specify the quality measure and how investment expenditure enters its assessment. Report coexistence of deterrence and productive coordination separately from the sign of \(W_1-W_0\). To attribute welfare dominance to a mechanism, declare the mechanism counterfactuals or decomposition rule and any interaction terms; the net welfare difference alone does not identify separate mechanism contributions.

For an empirical contribution, declare market units, horizon, observed and latent variables, measurement and sampling rules, an admissible model class \(\mathcal M\), its observable law \(P_O(m)\), and a target \(\tau(m)\) specifying the regime comparisons or mechanism contrasts under study. State the identifying variation and restrictions on latent contracts, selection and investment, and whether the result establishes point identification, informative bounds or a diagnostic implication, using the common identification conventions below. Contract terms and switching observations alone are not assumed sufficient. Claims of computational tractability must state their computational task and performance criterion; estimation claims must state their sampling assumptions and statistical criterion. These choices are application-specific requirements, not additional assumptions silently imposed on the entire agenda.
\subsection*{Variants and assumption changes}
\begin{enumerate}
\item \textbf{Editorial, fixed investment.} Compare exclusivity and nonexclusivity when entry cost is known and input quality fixed; existing exclusive-dealing examples are benchmarks, not a general unsolved assertion.
\item \textbf{Editorial, endogenous quality.} Permit investments before bargaining and rematching afterward; identify both exclusion and coordination channels from audited contract terms and investment data, with bounds when contracts are latent.
\end{enumerate}
\subsection*{References and source audit}
\textbf{Checked source.} 14 September 2021 author draft, Section 5, printed pp. 58--59. Full primary text retrieved. \url{https://robinlee.sites.fas.harvard.edu/papers/HandbookIO_Vol4_VerticalMarkets.pdf}.
\begin{enumerate}\raggedright
\item\hypertarget{definitions:source:30007681:1}{} Robin S. Lee; Michael D. Whinston; Ali Yurukoglu. 2021. \emph{Structural Empirical Analysis of Contracting in Vertical Markets}. 14 September 2021 author draft, Section 5, printed pp. 58--59.
\url{https://robinlee.sites.fas.harvard.edu/papers/HandbookIO_Vol4_VerticalMarkets.pdf}.
DOI: \url{https://doi.org/10.3386/w29282}.
\end{enumerate}

\subsection*{Common conventions: Source mathematical conventions}

These conventions apply to each entry unless explicitly replaced.

\textbf{Models and identification.} A primitive model \(m\) specifies all probability laws, feasible choices, information, equilibrium and measurement rules used by its target. The admissible class \(\mathcal M\) is fixed independently of outcomes. If \(P_O(m)\) is its observable probability law and \(\tau(m)\) the target, its identified set at observable law \(P\) is
\[
 \mathcal I_\tau(P)=\{\tau(m):m\in\mathcal M,\ P_O(m)=P\}.
\]
Point identification means that this set is a singleton. An empty set signals incompatibility of data and assumptions. Coordinate bounds need not describe the joint set. Bounds are sharp only if every retained target is attainable or arbitrarily approachable by compatible models. Finite bounds require explicit restrictions on unobserved quantities; unrestricted confounding, hidden wealth, extrapolation or arbitrary equilibrium selection can yield uninformative sets.

\textbf{Causal interventions.} A potential outcome \(Y_i(p)\) is the outcome under a complete policy regime \(p\), including timing, financing, information and market spillovers. The target population and units are fixed before treatment. With interference, use the complete assignment vector or a justified exposure mapping; individual no-interference notation alone cannot identify an economy-wide intervention. A difference of mean potential outcomes does not require the cross-world joint distribution. Conditioning on post-treatment migration, survival, enrollment or employment changes the estimand and needs separate justification.

\textbf{Statistical statements.} An estimator, confidence set, forecast or test specifies a sampling law, model class and loss/coverage criterion. Uniform \((1-\alpha)\)-coverage means \(\inf_{m\in\mathcal M}\Pr_m\{\tau(m)\in C_n\}\geq1-\alpha\), or an explicitly asymptotic version. The whole parameter space trivially has coverage, so informativeness requires a stated width, risk or power objective. A forecasting improvement is relative to a named benchmark, information set and evaluation population.

\textbf{Equilibrium and optimization.} An equilibrium comprises feasible plans, optimality, beliefs where needed, budget constraints and all market/resource clearing. The primitives or a stated benchmark define each correspondence \(\mathcal E(p;m)\). Nonemptiness is assumed or proved, never inferred just from compact policy sets. A selected-equilibrium target uses a declared selection rule; a robust target quantifies over all permitted equilibria. Use suprema or infima unless attainment is established. An \(\varepsilon\)-optimal policy is feasible and within \(\varepsilon\) of the finite optimum value. Report an empty feasible set separately.

\textbf{Welfare and accounting.} Normative utility, interpersonal normalization, social weights, numeraire, discounting and the use of public receipts are primitives. Behavioral utility need not coincide with the welfare criterion. Household welfare includes ownership income and rents once; adding profits again to compensating variation generally double-counts them. A monetary equivalent is defined by an explicit transfer and price convention, with monotonicity and range conditions. Average effects, mechanism contributions, transfers and welfare losses are different targets. Infinite sums require convergence and dynamic feasible plans require terminal or no-Ponzi conditions.

\textbf{Source versions.} A working-paper date, revision date and journal publication year are distinguished. A DOI for a journal version is not automatically the DOI of an earlier manuscript. Locators refer to the stated version and printed pages where specified. A metadata-only check does not verify a claim about a full-text conclusion. Every entry's source subsection narrows the provenance claim accordingly.
% END ECONOMICS STATEMENT
\end{document}
